\QHeading{第1题　在（\hspace{28pt}）里填上适当的数。}
\par\noindent


\PlainItem{$0.4$米$=$（\hspace{29pt}）分米{\color{answerred}$4$}$3.5$千米$=$（\hspace{27pt}）米}
\par

\Equation{\textcolor{answerred}{3500}}
\PlainItem{$0.02$吨$=$（\hspace{29pt}）千克{\color{answerred}$20$}$420$平方厘米$=$（\hspace{30pt}）平方分米}
\Equation{\textcolor{answerred}{4.2}}
\PlainItem{$5$米$7$分米$=$（\hspace{27pt}）米{\color{answerred}$5.7$}$5$元$7$分$=$（\hspace{28pt}）元}
\Equation{\textcolor{answerred}{5.07}}
\par\noindent

\Answer{
\QHeading{计算过程：第1题：复合单位换算}
5米7分米换算成米：
\Formula{
\begin{aligned}[t]&\phantom{={}}5+7\div10\\&=5+0.7\\&=5.7\text{（米）}\end{aligned}
}
5元7分换算成元：
\Formula{
\begin{aligned}[t]&\phantom{={}}5+7\div100\\&=5+0.07\\&=5.07\text{（元）}\end{aligned}
}

\QHeading{计算过程：第2题：列竖式计算的完整过程}
小数点对齐，缺位补零；从右往左逐列计算。每次进位、退位都在下面列出。
\par\medskip\noindent
\textbf{第1小题：\(9.00-6.28=2.72\)}\par\smallskip
\noindent
\begin{tikzpicture}[x=1em,y=1em,baseline=(current bounding box.center),draw=answerred,text=answerred]
\node[anchor=base,inner sep=0pt] at (0.000,0) {$9$};
\node[anchor=base,inner sep=0pt] at (0.800,0) {$.$};
\node[anchor=base,inner sep=0pt] at (1.100,0) {$0$};
\node[anchor=base,inner sep=0pt] at (1.900,0) {$0$};
\node[anchor=base,inner sep=0pt] at (0.000,-1.1) {$6$};
\node[anchor=base,inner sep=0pt] at (0.800,-1.1) {$.$};
\node[anchor=base,inner sep=0pt] at (1.100,-1.1) {$2$};
\node[anchor=base,inner sep=0pt] at (1.900,-1.1) {$8$};
\node[anchor=base,inner sep=0pt] at (0.000,-2.3) {$2$};
\node[anchor=base,inner sep=0pt] at (0.800,-2.3) {$.$};
\node[anchor=base,inner sep=0pt] at (1.100,-2.3) {$7$};
\node[anchor=base,inner sep=0pt] at (1.900,-2.3) {$2$};
\node[anchor=base,inner sep=0pt] at (-.7,-1.1) {$-$};
\draw[line width=.5pt] (-1,-1.52)--(2.500,-1.52);
\draw[line width=.35pt] (-0.160,.03)--(0.160,.43);
\node[anchor=base,inner sep=0pt,font=\normalsize] at (0.000,.72) {$8$};
\draw[line width=.35pt] (0.940,.03)--(1.260,.43);
\node[anchor=base,inner sep=0pt,font=\normalsize] at (1.100,.72) {$9$};
\draw[line width=.35pt] (1.740,.03)--(2.060,.43);
\node[anchor=base,inner sep=0pt,font=\normalsize] at (1.900,.72) {$10$};
\end{tikzpicture}
\par\smallskip{\normalsize 划线数字是退位前的数，上方小字是退位后参与计算的数。\par}
\vspace{0pt}
\begin{enumerate}[leftmargin=1.7em,itemsep=5pt,topsep=0pt]
\item 个位退位：\(\begin{aligned}[t]&\phantom{={}}9-1\\&=8\end{aligned}\)\quad 借出1给后一位。
\item 十分位传递退位：\(\begin{aligned}[t]&\phantom{={}}0+10-1\\&=10-1\\&=9\end{aligned}\)\quad 收到10，再借出1给后一位。
\item 百分位接收退位：\(\begin{aligned}[t]&\phantom{={}}0+10\\&=10\end{aligned}\)\quad 前一位的1变成本位的10。
\item 百分位：\(\begin{aligned}[t]&\phantom{={}}10-8\\&=2\end{aligned}\)\quad 本位写2。
\item 十分位：\(\begin{aligned}[t]&\phantom{={}}9-2\\&=7\end{aligned}\)\quad 本位写7。
\item 个位：\(\begin{aligned}[t]&\phantom{={}}8-6\\&=2\end{aligned}\)\quad 本位写2。
\end{enumerate}
\par\smallskip
因此，\(\begin{aligned}[t]&\phantom{={}}9.00-6.28\\&=2.72\end{aligned}\)。
\par\medskip

\par\medskip\noindent
\textbf{第2小题：\(1.70+1.85=3.55\)}\par\smallskip
\noindent
\begin{tikzpicture}[x=1em,y=1em,baseline=(current bounding box.center),draw=answerred,text=answerred]
\node[anchor=base,inner sep=0pt] at (0.000,0) {$1$};
\node[anchor=base,inner sep=0pt] at (0.800,0) {$.$};
\node[anchor=base,inner sep=0pt] at (1.100,0) {$7$};
\node[anchor=base,inner sep=0pt] at (1.900,0) {$0$};
\node[anchor=base,inner sep=0pt] at (0.000,-1.1) {$1$};
\node[anchor=base,inner sep=0pt] at (0.800,-1.1) {$.$};
\node[anchor=base,inner sep=0pt] at (1.100,-1.1) {$8$};
\node[anchor=base,inner sep=0pt] at (1.900,-1.1) {$5$};
\node[anchor=base,inner sep=0pt] at (0.000,-2.3) {$3$};
\node[anchor=base,inner sep=0pt] at (0.800,-2.3) {$.$};
\node[anchor=base,inner sep=0pt] at (1.100,-2.3) {$5$};
\node[anchor=base,inner sep=0pt] at (1.900,-2.3) {$5$};
\node[anchor=base,inner sep=0pt] at (-.7,-1.1) {$+$};
\draw[line width=.5pt] (-1,-1.52)--(2.500,-1.52);
\node[anchor=base,inner sep=0pt,font=\normalsize] at (0.000,.72) {$1$};
\end{tikzpicture}
\par\smallskip{\normalsize 上方小字表示从后一位进来的数。\par}
\vspace{0pt}
\begin{enumerate}[leftmargin=1.7em,itemsep=5pt,topsep=0pt]
\item 百分位：\(\begin{aligned}[t]&\phantom{={}}0+5\\&=5\end{aligned}\)\quad 本位写5。
\item 十分位：\(\begin{aligned}[t]&\phantom{={}}7+8\\&=15\end{aligned}\)\quad 本位写5，向前一位进1。
\item 个位：\(\begin{aligned}[t]&\phantom{={}}1+1+1\\&=2+1\\&=3\end{aligned}\)\quad 本位写3。
\end{enumerate}
\par\smallskip
因此，\(\begin{aligned}[t]&\phantom{={}}1.70+1.85\\&=3.55\end{aligned}\)。
\par\medskip

\par\medskip\noindent
\textbf{第3小题：\(2.85+1.05=3.9\)}\par\smallskip
\noindent
\begin{tikzpicture}[x=1em,y=1em,baseline=(current bounding box.center),draw=answerred,text=answerred]
\node[anchor=base,inner sep=0pt] at (0.000,0) {$2$};
\node[anchor=base,inner sep=0pt] at (0.800,0) {$.$};
\node[anchor=base,inner sep=0pt] at (1.100,0) {$8$};
\node[anchor=base,inner sep=0pt] at (1.900,0) {$5$};
\node[anchor=base,inner sep=0pt] at (0.000,-1.1) {$1$};
\node[anchor=base,inner sep=0pt] at (0.800,-1.1) {$.$};
\node[anchor=base,inner sep=0pt] at (1.100,-1.1) {$0$};
\node[anchor=base,inner sep=0pt] at (1.900,-1.1) {$5$};
\node[anchor=base,inner sep=0pt] at (0.000,-2.3) {$3$};
\node[anchor=base,inner sep=0pt] at (0.800,-2.3) {$.$};
\node[anchor=base,inner sep=0pt] at (1.100,-2.3) {$9$};
\node[anchor=base,inner sep=0pt] at (1.900,-2.3) {$0$};
\node[anchor=base,inner sep=0pt] at (-.7,-1.1) {$+$};
\draw[line width=.5pt] (-1,-1.52)--(2.500,-1.52);
\node[anchor=base,inner sep=0pt,font=\normalsize] at (1.100,.72) {$1$};
\end{tikzpicture}
\par\smallskip{\normalsize 上方小字表示从后一位进来的数。\par}
\vspace{0pt}
\begin{enumerate}[leftmargin=1.7em,itemsep=5pt,topsep=0pt]
\item 百分位：\(\begin{aligned}[t]&\phantom{={}}5+5\\&=10\end{aligned}\)\quad 本位写0，向前一位进1。
\item 十分位：\(\begin{aligned}[t]&\phantom{={}}8+0+1\\&=8+1\\&=9\end{aligned}\)\quad 本位写9。
\item 个位：\(\begin{aligned}[t]&\phantom{={}}2+1\\&=3\end{aligned}\)\quad 本位写3。
\end{enumerate}
\par\smallskip
因此，\(\begin{aligned}[t]&\phantom{={}}2.85+1.05\\&=3.90\end{aligned}\)。
\par\medskip

\par\medskip\noindent
\textbf{第4小题：\(5.90-3.58=2.32\)}\par\smallskip
\noindent
\begin{tikzpicture}[x=1em,y=1em,baseline=(current bounding box.center),draw=answerred,text=answerred]
\node[anchor=base,inner sep=0pt] at (0.000,0) {$5$};
\node[anchor=base,inner sep=0pt] at (0.800,0) {$.$};
\node[anchor=base,inner sep=0pt] at (1.100,0) {$9$};
\node[anchor=base,inner sep=0pt] at (1.900,0) {$0$};
\node[anchor=base,inner sep=0pt] at (0.000,-1.1) {$3$};
\node[anchor=base,inner sep=0pt] at (0.800,-1.1) {$.$};
\node[anchor=base,inner sep=0pt] at (1.100,-1.1) {$5$};
\node[anchor=base,inner sep=0pt] at (1.900,-1.1) {$8$};
\node[anchor=base,inner sep=0pt] at (0.000,-2.3) {$2$};
\node[anchor=base,inner sep=0pt] at (0.800,-2.3) {$.$};
\node[anchor=base,inner sep=0pt] at (1.100,-2.3) {$3$};
\node[anchor=base,inner sep=0pt] at (1.900,-2.3) {$2$};
\node[anchor=base,inner sep=0pt] at (-.7,-1.1) {$-$};
\draw[line width=.5pt] (-1,-1.52)--(2.500,-1.52);
\draw[line width=.35pt] (0.940,.03)--(1.260,.43);
\node[anchor=base,inner sep=0pt,font=\normalsize] at (1.100,.72) {$8$};
\draw[line width=.35pt] (1.740,.03)--(2.060,.43);
\node[anchor=base,inner sep=0pt,font=\normalsize] at (1.900,.72) {$10$};
\end{tikzpicture}
\par\smallskip{\normalsize 划线数字是退位前的数，上方小字是退位后参与计算的数。\par}
\vspace{0pt}
\begin{enumerate}[leftmargin=1.7em,itemsep=5pt,topsep=0pt]
\item 十分位退位：\(\begin{aligned}[t]&\phantom{={}}9-1\\&=8\end{aligned}\)\quad 借出1给后一位。
\item 百分位接收退位：\(\begin{aligned}[t]&\phantom{={}}0+10\\&=10\end{aligned}\)\quad 前一位的1变成本位的10。
\item 百分位：\(\begin{aligned}[t]&\phantom{={}}10-8\\&=2\end{aligned}\)\quad 本位写2。
\item 十分位：\(\begin{aligned}[t]&\phantom{={}}8-5\\&=3\end{aligned}\)\quad 本位写3。
\item 个位：\(\begin{aligned}[t]&\phantom{={}}5-3\\&=2\end{aligned}\)\quad 本位写2。
\end{enumerate}
\par\smallskip
因此，\(\begin{aligned}[t]&\phantom{={}}5.90-3.58\\&=2.32\end{aligned}\)。
\par\medskip

\par\medskip\noindent
\textbf{第5小题：\(2.52+1.03=3.55\)}\par\smallskip
\noindent
\begin{tikzpicture}[x=1em,y=1em,baseline=(current bounding box.center),draw=answerred,text=answerred]
\node[anchor=base,inner sep=0pt] at (0.000,0) {$2$};
\node[anchor=base,inner sep=0pt] at (0.800,0) {$.$};
\node[anchor=base,inner sep=0pt] at (1.100,0) {$5$};
\node[anchor=base,inner sep=0pt] at (1.900,0) {$2$};
\node[anchor=base,inner sep=0pt] at (0.000,-1.1) {$1$};
\node[anchor=base,inner sep=0pt] at (0.800,-1.1) {$.$};
\node[anchor=base,inner sep=0pt] at (1.100,-1.1) {$0$};
\node[anchor=base,inner sep=0pt] at (1.900,-1.1) {$3$};
\node[anchor=base,inner sep=0pt] at (0.000,-2.3) {$3$};
\node[anchor=base,inner sep=0pt] at (0.800,-2.3) {$.$};
\node[anchor=base,inner sep=0pt] at (1.100,-2.3) {$5$};
\node[anchor=base,inner sep=0pt] at (1.900,-2.3) {$5$};
\node[anchor=base,inner sep=0pt] at (-.7,-1.1) {$+$};
\draw[line width=.5pt] (-1,-1.52)--(2.500,-1.52);
\end{tikzpicture}
\par\smallskip{\normalsize 上方小字表示从后一位进来的数。\par}
\vspace{0pt}
\begin{enumerate}[leftmargin=1.7em,itemsep=5pt,topsep=0pt]
\item 百分位：\(\begin{aligned}[t]&\phantom{={}}2+3\\&=5\end{aligned}\)\quad 本位写5。
\item 十分位：\(\begin{aligned}[t]&\phantom{={}}5+0\\&=5\end{aligned}\)\quad 本位写5。
\item 个位：\(\begin{aligned}[t]&\phantom{={}}2+1\\&=3\end{aligned}\)\quad 本位写3。
\end{enumerate}
\par\smallskip
因此，\(\begin{aligned}[t]&\phantom{={}}2.52+1.03\\&=3.55\end{aligned}\)。
\par\medskip

\par\medskip\noindent
\textbf{第6小题：\(10.00-2.34=7.66\)}\par\smallskip
\noindent
\begin{tikzpicture}[x=1em,y=1em,baseline=(current bounding box.center),draw=answerred,text=answerred]
\node[anchor=base,inner sep=0pt] at (0.000,0) {$1$};
\node[anchor=base,inner sep=0pt] at (0.800,0) {$0$};
\node[anchor=base,inner sep=0pt] at (1.600,0) {$.$};
\node[anchor=base,inner sep=0pt] at (1.900,0) {$0$};
\node[anchor=base,inner sep=0pt] at (2.700,0) {$0$};
\node[anchor=base,inner sep=0pt] at (0.800,-1.1) {$2$};
\node[anchor=base,inner sep=0pt] at (1.600,-1.1) {$.$};
\node[anchor=base,inner sep=0pt] at (1.900,-1.1) {$3$};
\node[anchor=base,inner sep=0pt] at (2.700,-1.1) {$4$};
\node[anchor=base,inner sep=0pt] at (0.800,-2.3) {$7$};
\node[anchor=base,inner sep=0pt] at (1.600,-2.3) {$.$};
\node[anchor=base,inner sep=0pt] at (1.900,-2.3) {$6$};
\node[anchor=base,inner sep=0pt] at (2.700,-2.3) {$6$};
\node[anchor=base,inner sep=0pt] at (-.7,-1.1) {$-$};
\draw[line width=.5pt] (-1,-1.52)--(3.300,-1.52);
\draw[line width=.35pt] (-0.160,.03)--(0.160,.43);
\node[anchor=base,inner sep=0pt,font=\normalsize] at (0.000,.72) {$0$};
\draw[line width=.35pt] (0.640,.03)--(0.960,.43);
\node[anchor=base,inner sep=0pt,font=\normalsize] at (0.800,.72) {$9$};
\draw[line width=.35pt] (1.740,.03)--(2.060,.43);
\node[anchor=base,inner sep=0pt,font=\normalsize] at (1.900,.72) {$9$};
\draw[line width=.35pt] (2.540,.03)--(2.860,.43);
\node[anchor=base,inner sep=0pt,font=\normalsize] at (2.700,.72) {$10$};
\end{tikzpicture}
\par\smallskip{\normalsize 划线数字是退位前的数，上方小字是退位后参与计算的数。\par}
\vspace{0pt}
\begin{enumerate}[leftmargin=1.7em,itemsep=5pt,topsep=0pt]
\item 十位退位：\(\begin{aligned}[t]&\phantom{={}}1-1\\&=0\end{aligned}\)\quad 借出1给后一位。
\item 个位传递退位：\(\begin{aligned}[t]&\phantom{={}}0+10-1\\&=10-1\\&=9\end{aligned}\)\quad 收到10，再借出1给后一位。
\item 十分位传递退位：\(\begin{aligned}[t]&\phantom{={}}0+10-1\\&=10-1\\&=9\end{aligned}\)\quad 收到10，再借出1给后一位。
\item 百分位接收退位：\(\begin{aligned}[t]&\phantom{={}}0+10\\&=10\end{aligned}\)\quad 前一位的1变成本位的10。
\item 百分位：\(\begin{aligned}[t]&\phantom{={}}10-4\\&=6\end{aligned}\)\quad 本位写6。
\item 十分位：\(\begin{aligned}[t]&\phantom{={}}9-3\\&=6\end{aligned}\)\quad 本位写6。
\item 个位：\(\begin{aligned}[t]&\phantom{={}}9-2\\&=7\end{aligned}\)\quad 本位写7。
\item 十位：\(\begin{aligned}[t]&\phantom{={}}0-0\\&=0\end{aligned}\)\quad 本位写0。
\end{enumerate}
\par\smallskip
因此，\(\begin{aligned}[t]&\phantom{={}}10.00-2.34\\&=7.66\end{aligned}\)。
\par\medskip

\par\medskip\noindent
\textbf{第7小题：\(5.62-2.85=2.77\)}\par\smallskip
\noindent
\begin{tikzpicture}[x=1em,y=1em,baseline=(current bounding box.center),draw=answerred,text=answerred]
\node[anchor=base,inner sep=0pt] at (0.000,0) {$5$};
\node[anchor=base,inner sep=0pt] at (0.800,0) {$.$};
\node[anchor=base,inner sep=0pt] at (1.100,0) {$6$};
\node[anchor=base,inner sep=0pt] at (1.900,0) {$2$};
\node[anchor=base,inner sep=0pt] at (0.000,-1.1) {$2$};
\node[anchor=base,inner sep=0pt] at (0.800,-1.1) {$.$};
\node[anchor=base,inner sep=0pt] at (1.100,-1.1) {$8$};
\node[anchor=base,inner sep=0pt] at (1.900,-1.1) {$5$};
\node[anchor=base,inner sep=0pt] at (0.000,-2.3) {$2$};
\node[anchor=base,inner sep=0pt] at (0.800,-2.3) {$.$};
\node[anchor=base,inner sep=0pt] at (1.100,-2.3) {$7$};
\node[anchor=base,inner sep=0pt] at (1.900,-2.3) {$7$};
\node[anchor=base,inner sep=0pt] at (-.7,-1.1) {$-$};
\draw[line width=.5pt] (-1,-1.52)--(2.500,-1.52);
\draw[line width=.35pt] (-0.160,.03)--(0.160,.43);
\node[anchor=base,inner sep=0pt,font=\normalsize] at (0.000,.72) {$4$};
\draw[line width=.35pt] (0.940,.03)--(1.260,.43);
\node[anchor=base,inner sep=0pt,font=\normalsize] at (1.100,.72) {$15$};
\draw[line width=.35pt] (1.740,.03)--(2.060,.43);
\node[anchor=base,inner sep=0pt,font=\normalsize] at (1.900,.72) {$12$};
\end{tikzpicture}
\par\smallskip{\normalsize 划线数字是退位前的数，上方小字是退位后参与计算的数。\par}
\vspace{0pt}
\begin{enumerate}[leftmargin=1.7em,itemsep=5pt,topsep=0pt]
\item 十分位退位：\(\begin{aligned}[t]&\phantom{={}}6-1\\&=5\end{aligned}\)\quad 借出1给后一位。
\item 百分位接收退位：\(\begin{aligned}[t]&\phantom{={}}2+10\\&=12\end{aligned}\)\quad 前一位的1变成本位的10。
\item 百分位：\(\begin{aligned}[t]&\phantom{={}}12-5\\&=7\end{aligned}\)\quad 本位写7。
\item 个位退位：\(\begin{aligned}[t]&\phantom{={}}5-1\\&=4\end{aligned}\)\quad 借出1给后一位。
\item 十分位接收退位：\(\begin{aligned}[t]&\phantom{={}}5+10\\&=15\end{aligned}\)\quad 前一位的1变成本位的10。
\item 十分位：\(\begin{aligned}[t]&\phantom{={}}15-8\\&=7\end{aligned}\)\quad 本位写7。
\item 个位：\(\begin{aligned}[t]&\phantom{={}}4-2\\&=2\end{aligned}\)\quad 本位写2。
\end{enumerate}
\par\smallskip
因此，\(\begin{aligned}[t]&\phantom{={}}5.62-2.85\\&=2.77\end{aligned}\)。
\par\medskip

\par\medskip\noindent
\textbf{第8小题：\(3.27+5.79=9.06\)}\par\smallskip
\noindent
\begin{tikzpicture}[x=1em,y=1em,baseline=(current bounding box.center),draw=answerred,text=answerred]
\node[anchor=base,inner sep=0pt] at (0.000,0) {$3$};
\node[anchor=base,inner sep=0pt] at (0.800,0) {$.$};
\node[anchor=base,inner sep=0pt] at (1.100,0) {$2$};
\node[anchor=base,inner sep=0pt] at (1.900,0) {$7$};
\node[anchor=base,inner sep=0pt] at (0.000,-1.1) {$5$};
\node[anchor=base,inner sep=0pt] at (0.800,-1.1) {$.$};
\node[anchor=base,inner sep=0pt] at (1.100,-1.1) {$7$};
\node[anchor=base,inner sep=0pt] at (1.900,-1.1) {$9$};
\node[anchor=base,inner sep=0pt] at (0.000,-2.3) {$9$};
\node[anchor=base,inner sep=0pt] at (0.800,-2.3) {$.$};
\node[anchor=base,inner sep=0pt] at (1.100,-2.3) {$0$};
\node[anchor=base,inner sep=0pt] at (1.900,-2.3) {$6$};
\node[anchor=base,inner sep=0pt] at (-.7,-1.1) {$+$};
\draw[line width=.5pt] (-1,-1.52)--(2.500,-1.52);
\node[anchor=base,inner sep=0pt,font=\normalsize] at (1.100,.72) {$1$};
\node[anchor=base,inner sep=0pt,font=\normalsize] at (0.000,.72) {$1$};
\end{tikzpicture}
\par\smallskip{\normalsize 上方小字表示从后一位进来的数。\par}
\vspace{0pt}
\begin{enumerate}[leftmargin=1.7em,itemsep=5pt,topsep=0pt]
\item 百分位：\(\begin{aligned}[t]&\phantom{={}}7+9\\&=16\end{aligned}\)\quad 本位写6，向前一位进1。
\item 十分位：\(\begin{aligned}[t]&\phantom{={}}2+7+1\\&=9+1\\&=10\end{aligned}\)\quad 本位写0，向前一位进1。
\item 个位：\(\begin{aligned}[t]&\phantom{={}}3+5+1\\&=8+1\\&=9\end{aligned}\)\quad 本位写9。
\end{enumerate}
\par\smallskip
因此，\(\begin{aligned}[t]&\phantom{={}}3.27+5.79\\&=9.06\end{aligned}\)。
\par\medskip

}

\QHeading{第2题　列竖式计算。}

\Problem{（1）$9-6.28=2.72$}
\par


\MethodTitle{方法一：按题意列竖式}
\Answer{\begin{tikzpicture}[x=1em,y=1em,baseline=(current bounding box.center),draw=answerred,text=answerred]
\node[anchor=base,inner sep=0pt] at (0.000,0) {$9$};
\node[anchor=base,inner sep=0pt] at (0.800,0) {$.$};
\node[anchor=base,inner sep=0pt] at (1.120,0) {$0$};
\node[anchor=base,inner sep=0pt] at (1.920,0) {$0$};
\node[anchor=base,inner sep=0pt] at (0.000,-1.1) {$6$};
\node[anchor=base,inner sep=0pt] at (0.800,-1.1) {$.$};
\node[anchor=base,inner sep=0pt] at (1.120,-1.1) {$2$};
\node[anchor=base,inner sep=0pt] at (1.920,-1.1) {$8$};
\node[anchor=base,inner sep=0pt] at (0.000,-2.3) {$2$};
\node[anchor=base,inner sep=0pt] at (0.800,-2.3) {$.$};
\node[anchor=base,inner sep=0pt] at (1.120,-2.3) {$7$};
\node[anchor=base,inner sep=0pt] at (1.920,-2.3) {$2$};
\node[anchor=base,inner sep=0pt] at (-.8,-1.1) {$-$};
\draw[line width=.5pt] (-1.1,-1.52)--(2.440,-1.52);
\draw[line width=.35pt] (-0.190,.03)--(0.190,.47);
\node[anchor=base,inner sep=0pt,font=\normalsize] at (0.000,.77) {$8$};
\draw[line width=.35pt] (0.930,.03)--(1.310,.47);
\node[anchor=base,inner sep=0pt,font=\normalsize] at (1.120,.77) {$9$};
\draw[line width=.35pt] (1.730,.03)--(2.110,.47);
\node[anchor=base,inner sep=0pt,font=\normalsize] at (1.920,.77) {$10$};
\end{tikzpicture}}

\MethodBlock{方法二：拆成整数部分和小数部分}{上方竖式及题目要求的验算仍保留；本法用于核对。}{\begin{aligned}[t]&\phantom{={}}9-6.28\\
&=9-6-0.28\\
&=3-0.28\\
&=2.72\end{aligned}}

\MethodBlock{方法三：按相同计数单位计算}{上方竖式及题目要求的验算仍保留；本法用于核对。先看成若干个百分之一，按整数加减，最后还原成小数。}{\begin{aligned}[t]&\phantom{={}}9-6.28\\
&=(900-628)\div 100\\
&=272\div 100\\
&=2.72\end{aligned}}
\Problem{（2）$1.7+1.85=3.55$}

\MethodTitle{方法一：按题意列竖式}
\Answer{\begin{tikzpicture}[x=1em,y=1em,baseline=(current bounding box.center),draw=answerred,text=answerred]
\node[anchor=base,inner sep=0pt] at (0.000,0) {$1$};
\node[anchor=base,inner sep=0pt] at (0.800,0) {$.$};
\node[anchor=base,inner sep=0pt] at (1.120,0) {$7$};
\node[anchor=base,inner sep=0pt] at (1.920,0) {$0$};
\node[anchor=base,inner sep=0pt] at (0.000,-1.1) {$1$};
\node[anchor=base,inner sep=0pt] at (0.800,-1.1) {$.$};
\node[anchor=base,inner sep=0pt] at (1.120,-1.1) {$8$};
\node[anchor=base,inner sep=0pt] at (1.920,-1.1) {$5$};
\node[anchor=base,inner sep=0pt] at (0.000,-2.3) {$3$};
\node[anchor=base,inner sep=0pt] at (0.800,-2.3) {$.$};
\node[anchor=base,inner sep=0pt] at (1.120,-2.3) {$5$};
\node[anchor=base,inner sep=0pt] at (1.920,-2.3) {$5$};
\node[anchor=base,inner sep=0pt] at (-.8,-1.1) {$+$};
\draw[line width=.5pt] (-1.1,-1.52)--(2.440,-1.52);
\node[anchor=base,inner sep=0pt,font=\normalsize] at (0.000,.77) {$1$};
\end{tikzpicture}}

\MethodBlock{方法二：凑整再减去多加的数}{上方竖式及题目要求的验算仍保留；本法用于核对。把第二个数凑成整数，再减去多加的部分。}{\begin{aligned}[t]&\phantom{={}}1.7+1.85\\
&=1.7+2-0.15\\
&=3.7-0.15\\
&=3.55\end{aligned}}

\MethodBlock{方法三：拆成整数部分和小数部分}{上方竖式及题目要求的验算仍保留；本法用于核对。}{\begin{aligned}[t]&\phantom{={}}1.7+1.85\\
&=1.7+1+0.85\\
&=2.7+0.85\\
&=3.55\end{aligned}}
\Problem{（3）$2.85+1.05=3.9$}

\MethodTitle{方法一：按题意列竖式}
\Answer{\begin{tikzpicture}[x=1em,y=1em,baseline=(current bounding box.center),draw=answerred,text=answerred]
\node[anchor=base,inner sep=0pt] at (0.000,0) {$2$};
\node[anchor=base,inner sep=0pt] at (0.800,0) {$.$};
\node[anchor=base,inner sep=0pt] at (1.120,0) {$8$};
\node[anchor=base,inner sep=0pt] at (1.920,0) {$5$};
\node[anchor=base,inner sep=0pt] at (0.000,-1.1) {$1$};
\node[anchor=base,inner sep=0pt] at (0.800,-1.1) {$.$};
\node[anchor=base,inner sep=0pt] at (1.120,-1.1) {$0$};
\node[anchor=base,inner sep=0pt] at (1.920,-1.1) {$5$};
\node[anchor=base,inner sep=0pt] at (0.000,-2.3) {$3$};
\node[anchor=base,inner sep=0pt] at (0.800,-2.3) {$.$};
\node[anchor=base,inner sep=0pt] at (1.120,-2.3) {$9$};
\node[anchor=base,inner sep=0pt] at (1.920,-2.3) {$0$};
\node[anchor=base,inner sep=0pt] at (-.8,-1.1) {$+$};
\draw[line width=.5pt] (-1.1,-1.52)--(2.440,-1.52);
\node[anchor=base,inner sep=0pt,font=\normalsize] at (1.120,.77) {$1$};
\end{tikzpicture}}

\MethodBlock{方法二：拆成整数部分和小数部分}{上方竖式及题目要求的验算仍保留；本法用于核对。}{\begin{aligned}[t]&\phantom{={}}2.85+1.05\\
&=2.85+1+0.05\\
&=3.85+0.05\\
&=3.9\end{aligned}}

\MethodBlock{方法三：按相同计数单位计算}{上方竖式及题目要求的验算仍保留；本法用于核对。先看成若干个百分之一，按整数加减，最后还原成小数。}{\begin{aligned}[t]&\phantom{={}}2.85+1.05\\
&=(285+105)\div 100\\
&=390\div 100\\
&=3.9\end{aligned}}
\Problem{（4）$5.9-3.58=2.32$}

\MethodTitle{方法一：按题意列竖式}
\Answer{\begin{tikzpicture}[x=1em,y=1em,baseline=(current bounding box.center),draw=answerred,text=answerred]
\node[anchor=base,inner sep=0pt] at (0.000,0) {$5$};
\node[anchor=base,inner sep=0pt] at (0.800,0) {$.$};
\node[anchor=base,inner sep=0pt] at (1.120,0) {$9$};
\node[anchor=base,inner sep=0pt] at (1.920,0) {$0$};
\node[anchor=base,inner sep=0pt] at (0.000,-1.1) {$3$};
\node[anchor=base,inner sep=0pt] at (0.800,-1.1) {$.$};
\node[anchor=base,inner sep=0pt] at (1.120,-1.1) {$5$};
\node[anchor=base,inner sep=0pt] at (1.920,-1.1) {$8$};
\node[anchor=base,inner sep=0pt] at (0.000,-2.3) {$2$};
\node[anchor=base,inner sep=0pt] at (0.800,-2.3) {$.$};
\node[anchor=base,inner sep=0pt] at (1.120,-2.3) {$3$};
\node[anchor=base,inner sep=0pt] at (1.920,-2.3) {$2$};
\node[anchor=base,inner sep=0pt] at (-.8,-1.1) {$-$};
\draw[line width=.5pt] (-1.1,-1.52)--(2.440,-1.52);
\draw[line width=.35pt] (0.930,.03)--(1.310,.47);
\node[anchor=base,inner sep=0pt,font=\normalsize] at (1.120,.77) {$8$};
\draw[line width=.35pt] (1.730,.03)--(2.110,.47);
\node[anchor=base,inner sep=0pt,font=\normalsize] at (1.920,.77) {$10$};
\end{tikzpicture}}

\MethodBlock{方法二：拆成整数部分和小数部分}{上方竖式及题目要求的验算仍保留；本法用于核对。}{\begin{aligned}[t]&\phantom{={}}5.9-3.58\\
&=5.9-3-0.58\\
&=2.9-0.58\\
&=2.32\end{aligned}}

\MethodBlock{方法三：按相同计数单位计算}{上方竖式及题目要求的验算仍保留；本法用于核对。先看成若干个百分之一，按整数加减，最后还原成小数。}{\begin{aligned}[t]&\phantom{={}}5.9-3.58\\
&=(590-358)\div 100\\
&=232\div 100\\
&=2.32\end{aligned}}
\Problem{（5）$2.52+1.03=3.55$}

\MethodTitle{方法一：按题意列竖式}
\Answer{\begin{tikzpicture}[x=1em,y=1em,baseline=(current bounding box.center),draw=answerred,text=answerred]
\node[anchor=base,inner sep=0pt] at (0.000,0) {$2$};
\node[anchor=base,inner sep=0pt] at (0.800,0) {$.$};
\node[anchor=base,inner sep=0pt] at (1.120,0) {$5$};
\node[anchor=base,inner sep=0pt] at (1.920,0) {$2$};
\node[anchor=base,inner sep=0pt] at (0.000,-1.1) {$1$};
\node[anchor=base,inner sep=0pt] at (0.800,-1.1) {$.$};
\node[anchor=base,inner sep=0pt] at (1.120,-1.1) {$0$};
\node[anchor=base,inner sep=0pt] at (1.920,-1.1) {$3$};
\node[anchor=base,inner sep=0pt] at (0.000,-2.3) {$3$};
\node[anchor=base,inner sep=0pt] at (0.800,-2.3) {$.$};
\node[anchor=base,inner sep=0pt] at (1.120,-2.3) {$5$};
\node[anchor=base,inner sep=0pt] at (1.920,-2.3) {$5$};
\node[anchor=base,inner sep=0pt] at (-.8,-1.1) {$+$};
\draw[line width=.5pt] (-1.1,-1.52)--(2.440,-1.52);
\end{tikzpicture}}

\MethodBlock{方法二：拆成整数部分和小数部分}{上方竖式及题目要求的验算仍保留；本法用于核对。}{\begin{aligned}[t]&\phantom{={}}2.52+1.03\\
&=2.52+1+0.03\\
&=3.52+0.03\\
&=3.55\end{aligned}}

\MethodBlock{方法三：按相同计数单位计算}{上方竖式及题目要求的验算仍保留；本法用于核对。先看成若干个百分之一，按整数加减，最后还原成小数。}{\begin{aligned}[t]&\phantom{={}}2.52+1.03\\
&=(252+103)\div 100\\
&=355\div 100\\
&=3.55\end{aligned}}
\Problem{（6）$10-2.34=7.66$}

\MethodTitle{方法一：按题意列竖式}
\Answer{\begin{tikzpicture}[x=1em,y=1em,baseline=(current bounding box.center),draw=answerred,text=answerred]
\node[anchor=base,inner sep=0pt] at (0.000,0) {$1$};
\node[anchor=base,inner sep=0pt] at (0.800,0) {$0$};
\node[anchor=base,inner sep=0pt] at (1.600,0) {$.$};
\node[anchor=base,inner sep=0pt] at (1.920,0) {$0$};
\node[anchor=base,inner sep=0pt] at (2.720,0) {$0$};
\node[anchor=base,inner sep=0pt] at (0.800,-1.1) {$2$};
\node[anchor=base,inner sep=0pt] at (1.600,-1.1) {$.$};
\node[anchor=base,inner sep=0pt] at (1.920,-1.1) {$3$};
\node[anchor=base,inner sep=0pt] at (2.720,-1.1) {$4$};
\node[anchor=base,inner sep=0pt] at (0.800,-2.3) {$7$};
\node[anchor=base,inner sep=0pt] at (1.600,-2.3) {$.$};
\node[anchor=base,inner sep=0pt] at (1.920,-2.3) {$6$};
\node[anchor=base,inner sep=0pt] at (2.720,-2.3) {$6$};
\node[anchor=base,inner sep=0pt] at (-.8,-1.1) {$-$};
\draw[line width=.5pt] (-1.1,-1.52)--(3.240,-1.52);
\draw[line width=.35pt] (-0.190,.03)--(0.190,.47);
\node[anchor=base,inner sep=0pt,font=\normalsize] at (0.000,.77) {$0$};
\draw[line width=.35pt] (0.610,.03)--(0.990,.47);
\node[anchor=base,inner sep=0pt,font=\normalsize] at (0.800,.77) {$9$};
\draw[line width=.35pt] (1.730,.03)--(2.110,.47);
\node[anchor=base,inner sep=0pt,font=\normalsize] at (1.920,.77) {$9$};
\draw[line width=.35pt] (2.530,.03)--(2.910,.47);
\node[anchor=base,inner sep=0pt,font=\normalsize] at (2.720,.77) {$10$};
\end{tikzpicture}}

\MethodBlock{方法二：拆成整数部分和小数部分}{上方竖式及题目要求的验算仍保留；本法用于核对。}{\begin{aligned}[t]&\phantom{={}}10-2.34\\
&=10-2-0.34\\
&=8-0.34\\
&=7.66\end{aligned}}

\MethodBlock{方法三：按相同计数单位计算}{上方竖式及题目要求的验算仍保留；本法用于核对。先看成若干个百分之一，按整数加减，最后还原成小数。}{\begin{aligned}[t]&\phantom{={}}10-2.34\\
&=(1000-234)\div 100\\
&=766\div 100\\
&=7.66\end{aligned}}
\Problem{（7）$5.62-2.85=2.77$}

\MethodTitle{方法一：按题意列竖式}
\Answer{\begin{tikzpicture}[x=1em,y=1em,baseline=(current bounding box.center),draw=answerred,text=answerred]
\node[anchor=base,inner sep=0pt] at (0.000,0) {$5$};
\node[anchor=base,inner sep=0pt] at (0.800,0) {$.$};
\node[anchor=base,inner sep=0pt] at (1.120,0) {$6$};
\node[anchor=base,inner sep=0pt] at (1.920,0) {$2$};
\node[anchor=base,inner sep=0pt] at (0.000,-1.1) {$2$};
\node[anchor=base,inner sep=0pt] at (0.800,-1.1) {$.$};
\node[anchor=base,inner sep=0pt] at (1.120,-1.1) {$8$};
\node[anchor=base,inner sep=0pt] at (1.920,-1.1) {$5$};
\node[anchor=base,inner sep=0pt] at (0.000,-2.3) {$2$};
\node[anchor=base,inner sep=0pt] at (0.800,-2.3) {$.$};
\node[anchor=base,inner sep=0pt] at (1.120,-2.3) {$7$};
\node[anchor=base,inner sep=0pt] at (1.920,-2.3) {$7$};
\node[anchor=base,inner sep=0pt] at (-.8,-1.1) {$-$};
\draw[line width=.5pt] (-1.1,-1.52)--(2.440,-1.52);
\draw[line width=.35pt] (-0.190,.03)--(0.190,.47);
\node[anchor=base,inner sep=0pt,font=\normalsize] at (0.000,.77) {$4$};
\draw[line width=.35pt] (0.930,.03)--(1.310,.47);
\node[anchor=base,inner sep=0pt,font=\normalsize] at (1.120,.77) {$15$};
\draw[line width=.35pt] (1.730,.03)--(2.110,.47);
\node[anchor=base,inner sep=0pt,font=\normalsize] at (1.920,.77) {$12$};
\end{tikzpicture}}

\MethodBlock{方法二：多减再加回}{上方竖式及题目要求的验算仍保留；本法用于核对。减去凑成的整数后，把多减的部分加回来。}{\begin{aligned}[t]&\phantom{={}}5.62-2.85\\
&=5.62-3+0.15\\
&=2.62+0.15\\
&=2.77\end{aligned}}

\MethodBlock{方法三：拆成整数部分和小数部分}{上方竖式及题目要求的验算仍保留；本法用于核对。}{\begin{aligned}[t]&\phantom{={}}5.62-2.85\\
&=5.62-2-0.85\\
&=3.62-0.85\\
&=2.77\end{aligned}}
\Problem{（8）$3.27+5.79=9.06$}

\MethodTitle{方法一：按题意列竖式}
\Answer{\begin{tikzpicture}[x=1em,y=1em,baseline=(current bounding box.center),draw=answerred,text=answerred]
\node[anchor=base,inner sep=0pt] at (0.000,0) {$3$};
\node[anchor=base,inner sep=0pt] at (0.800,0) {$.$};
\node[anchor=base,inner sep=0pt] at (1.120,0) {$2$};
\node[anchor=base,inner sep=0pt] at (1.920,0) {$7$};
\node[anchor=base,inner sep=0pt] at (0.000,-1.1) {$5$};
\node[anchor=base,inner sep=0pt] at (0.800,-1.1) {$.$};
\node[anchor=base,inner sep=0pt] at (1.120,-1.1) {$7$};
\node[anchor=base,inner sep=0pt] at (1.920,-1.1) {$9$};
\node[anchor=base,inner sep=0pt] at (0.000,-2.3) {$9$};
\node[anchor=base,inner sep=0pt] at (0.800,-2.3) {$.$};
\node[anchor=base,inner sep=0pt] at (1.120,-2.3) {$0$};
\node[anchor=base,inner sep=0pt] at (1.920,-2.3) {$6$};
\node[anchor=base,inner sep=0pt] at (-.8,-1.1) {$+$};
\draw[line width=.5pt] (-1.1,-1.52)--(2.440,-1.52);
\node[anchor=base,inner sep=0pt,font=\normalsize] at (1.120,.77) {$1$};
\node[anchor=base,inner sep=0pt,font=\normalsize] at (0.000,.77) {$1$};
\end{tikzpicture}}

\MethodBlock{方法二：凑整再减去多加的数}{上方竖式及题目要求的验算仍保留；本法用于核对。把第二个数凑成整数，再减去多加的部分。}{\begin{aligned}[t]&\phantom{={}}3.27+5.79\\
&=3.27+6-0.21\\
&=9.27-0.21\\
&=9.06\end{aligned}}

\MethodBlock{方法三：拆成整数部分和小数部分}{上方竖式及题目要求的验算仍保留；本法用于核对。}{\begin{aligned}[t]&\phantom{={}}3.27+5.79\\
&=3.27+5+0.79\\
&=8.27+0.79\\
&=9.06\end{aligned}}
\par\noindent

\QHeading{第3题　解决问题。}

\Problem{① 买一个足球需要77.8元，一个足球比一个排球便宜2.12元。买一个足球和一个排球一共需要多少钱？}
\par


\MethodTitle{方法一}
\Answer{{\Formula{\begin{aligned}[t]&\phantom{={}}77.8+(77.8+2.12)\\&=77.8+79.92\\&=157.72\text{（元）}\end{aligned}}}\par 答：一共需要157.72元。}

\MethodTitle{方法二：先按两球同价计算，再补差价}
\Answer{排球比足球贵2.12元，两个都先按足球价77.8元算，再加差额。}
\Answer{\Formula{\begin{aligned}[t]&\phantom{={}}77.8\times 2+2.12\\
&=155.6+2.12\\
&=157.72\text{（元）}\end{aligned}}}
\Answer{答：一共需要157.72元。}
\Problem{② 四年级学生参加元旦文艺演出，跳舞的有34人，合唱的有50人，两项都参加的有13人。一共有多少人参加演出？}

\MethodTitle{方法一}
\Answer{{\Formula{\begin{aligned}[t]&\phantom{={}}34+50-13\\&=84-13\\&=71\text{（人）}\end{aligned}}}\par 答：一共有71人参加演出。}


\MethodTitle{方法二：先求只属于第二类的人数}
\Answer{第一类人数已经包含两类兼有的人，只需再加上仅属于第二类的人。}
\Answer{\Formula{\begin{aligned}[t]&\phantom{={}}34+(50-13)\\
&=34+37\\
&=71\text{（人）}\end{aligned}}}
\Answer{答：参加演出的一共有71人。}

\MethodTitle{方法三：分成互不重叠的三类}
\Answer{分别计数“只属于第一类”“两类兼有”“只属于第二类”，每人只数一次。}
\Answer{\Formula{\begin{aligned}[t]&\phantom{={}}(34-13)+13+(50-13)\\
&=21+13+37\\
&=34+37\\
&=71\text{（人）}\end{aligned}}}
\Answer{答：参加演出的一共有71人。}

